

\RequirePackage{fix-cm}
\documentclass[letterpaper,13pt]{article}

\usepackage{latexsym}
\usepackage[empty]{fullpage}
\usepackage{titlesec}
\usepackage{marvosym}
\usepackage[usenames,dvipsnames]{color}
\usepackage{verbatim}
\usepackage{enumitem}
\usepackage[hidelinks]{hyperref}
\usepackage{fancyhdr}
\usepackage[english]{babel}
\usepackage{tabularx}
\usepackage[absolute,overlay]{textpos}
\usepackage{xcolor}
\input{glyphtounicode}


\pagestyle{fancy}
\fancyhf{} 
\fancyfoot{}
\renewcommand{\headrulewidth}{0pt}
\renewcommand{\footrulewidth}{0pt}

\addtolength{\oddsidemargin}{-0.5in}
\addtolength{\evensidemargin}{-0.5in}
\addtolength{\textwidth}{1in}
\addtolength{\topmargin}{-.55in}
\addtolength{\textheight}{1.15in}

\urlstyle{same}

\raggedbottom
\raggedright
\setlength{\tabcolsep}{0in}

\titleformat{\section}{
  \vspace{0pt}\scshape\raggedright\large
}{}{0em}{}[\color{black}\titlerule \vspace{-2pt}]

\pdfgentounicode=1


\newcommand{\resumeItem}[1]{
  \item\small{
    {#1}
  }
}

\newcommand{\resumeSubheading}[4]{
  \vspace{-2pt}\item
    \begin{tabular*}{\textwidth}[t]{l@{\extracolsep{\fill}}r}
      \textbf{#1} & #2 \enspace $|$ \enspace \textit{\small #4} \\
      {\fontsize{10pt}{11pt}\selectfont\textit{#3}} & \\
    \end{tabular*}\vspace{-1pt}
}

\newcommand{\resumeSubSubheading}[2]{
    \item
    \begin{tabular*}{\textwidth}{l@{\extracolsep{\fill}}r}
      \textit{\small#1} & \textit{\small #2} \\
    \end{tabular*}\vspace{-5pt}
}

\newcommand{\resumeProjectHeading}[2]{
    \item
    \begin{tabular*}{\textwidth}{l@{\extracolsep{\fill}}r}
      \small#1 & #2 \\
    \end{tabular*}\vspace{-4pt}
}

\newcommand{\resumeSubItem}[1]{\resumeItem{#1}\vspace{-4pt}}

\renewcommand\labelitemii{$\vcenter{\hbox{\tiny$\bullet$}}$}

\newcommand{\resumeSubHeadingListStart}{\begin{itemize}[leftmargin=0in, label={}]}
\newcommand{\resumeSubHeadingListEnd}{\end{itemize}}
\newcommand{\resumeItemListStart}{\begin{itemize}[leftmargin=0.22in, itemsep=2pt, parsep=0pt, topsep=1pt, partopsep=0pt]}
\newcommand{\resumeItemListEnd}{\end{itemize}\vspace{-2pt}}

\begin{document}

\begin{center}
    \textbf{\LARGE \scshape Mohammadreza Nemati} \\ \vspace{1pt}
    \small  Santa Clara, CA 95054 $|$ 
    \small Mobile: 567-204-7520 $|$ 
    \href{mailto:mrnemati747@gmail.com}{mrnemati747@gmail.com} $|$
    \small \href{https://www.linkedin.com/in/mnemati/}{\textcolor{blue}{\underline{www.linkedin.com/in/mnemati}}}
\end{center}

\section{Summary}
\small{AI Engineer and Computer Science PhD specializing in agentic AI, multi-agent frameworks, retrieval-augmented generation, and production LLM systems. Experienced building and evaluating autonomous agents, fine-tuning foundation models, and developing scalable pipelines for debugging, deployment, and workflow automation.}
\vspace{0.05cm}

\section{Experience}
\vspace{0.05cm}
\resumeSubHeadingListStart

\resumeSubheading
  {AI Engineer}{2025 - Present}
  {Volkswagen}{Belmont, CA}
  \resumeItemListStart
    \resumeItem{\textbf{Multi-Agent Frameworks:} Built a production-grade LangGraph-based multimodal multi-agent system for long-horizon Android and browser workflows, orchestrating task decomposition, tool execution, verification, dynamic replanning, and error recovery on real devices.}
    \resumeItem{\textbf{Retrieval-Augmented Generation:} Integrated RAG into LangGraph agents to retrieve relevant task context and execution histories for grounded planning, verification, troubleshooting, and error recovery.}
    \resumeItem{\textbf{Agent Evaluation \& Debugging:} Engineered evaluation and observability infrastructure that logged screenshots, model outputs, tool actions, latency, execution plans, and step traces to isolate failure modes and measure end-to-end task success.}
    \resumeItem{\textbf{Foundation-Model Fine-Tuning:} Post-trained multimodal foundation models with SFT, DPO, and GRPO-based reinforcement learning using distributed PyTorch pipelines with DDP for planning, reasoning, grounding, verification, and error recovery.}
    \resumeItem{\textbf{Production LLM Deployment:} Deployed and evaluated Qwen3.5 2B with vLLM, tuning prefix caching, long-context inference, and serving parameters to improve maximum request concurrency.}
    \resumeItem{\textbf{Training Data Pipelines:} Continually pre-trained Qwen3.5 4B for use as a GUI-agent world model on approximately 500,000 trajectories assembled from public benchmarks and custom in-house data, averaging 22 interaction turns per trajectory.}
  \resumeItemListEnd
  \vspace{0.04cm}

\resumeSubheading
  {AI Researcher}{2022 - 2025}
  {Case Western Reserve University}{Cleveland, Ohio}
  \resumeItemListStart
    \resumeItem{\textbf{LLM Inference Optimization:} Developed an adaptive mixed-precision KV-cache quantization method based on key/value norm imbalance, achieving 2$\times$ memory reduction with $>$94\% accuracy retention across LLMs up to 70B parameters.}
    \resumeItem{\textbf{Model Development:} Developed LIT-LVM, a latent-variable method for structured interaction modeling that outperformed elastic net, hierarchical lasso, and factorization-machine baselines across simulated and real-world datasets. Published in Transactions on Machine Learning Research (2025).}
    \resumeItem{\textbf{Model Development:} Developed DEASurv, a multimodal time-to-event survival prediction model combining entity-specific embeddings, attention, and contrastive learning for large-scale medical datasets.}
  \resumeItemListEnd
  \vspace{0.02cm}

\resumeSubheading
  {Software Engineer - AI/ML Intern}{2024}
  {Sherwin-Williams}{Cleveland, Ohio}
  \resumeItemListStart
    \resumeItem{\textbf{Perception Model Development:} Built and deployed a VLM-based perception system integrating LLaVA, vision transformers, and SAM for natural-language object detection, grounding, and semantic segmentation, improving segmentation quality by 7\% IoU.}
  \resumeItemListEnd
  \vspace{0.02cm}

\resumeSubheading
  {Data Science Research Intern}{2023}
  {Beyond AI}{Glendale, California}
  \resumeItemListStart
    \resumeItem{\textbf{LLM Fine-Tuning:} Fine-tuned Llama-2 on refinery-domain text using PyTorch FSDP/DDP for scalable multi-GPU training and quantized deployment.}
  \resumeItemListEnd
  \vspace{0.02cm}

\resumeSubheading
  {Software Intern}{2020}
  {Nexteer Automotive}{Auburn Hills, Michigan}
  \resumeItemListStart
    \resumeItem{\textbf{Data Pipeline:} Developed a data pipeline for electric power-steering test data, enabling structured analysis and validation of steering-system responses for ADAS applications.}
  \resumeItemListEnd
\resumeSubHeadingListEnd

\section{Skills}
\vspace{0.02cm}
\begin{itemize}[leftmargin=0.15in, label={}, itemsep=2pt, topsep=1pt]
  \item \textbf{Programming}: Python, C++, SQL
  \item \textbf{Agentic AI}: LangGraph, Multi-Agent Systems, Retrieval-Augmented Generation, Tool Use, Dynamic Planning
  \item \textbf{Machine Learning}: PyTorch, Transformers, Large Language Models, Reinforcement Learning, Model Evaluation
  \item \textbf{ML Systems}: vLLM, Agent Observability, Distributed Training, FSDP, DDP, Docker
\end{itemize}
\vspace{0.02cm}

\section{Education}
\resumeSubHeadingListStart
  \resumeProjectHeading{\textbf{Case Western Reserve University} -- PhD, Computer Science}{Cleveland, OH $|$ 2025}
\resumeSubHeadingListEnd
\end{document}
